\documentclass{iopjournal}
\usepackage{amssymb,amsmath,booktabs,array,multirow,adjustbox,makecell,microtype}
\usepackage[round]{natbib}
\usepackage{placeins,ragged2e}
\hypersetup{hidelinks,pdftitle={Subject-independent forearm sEMG electrode reduction with cross-fitted attribution and non-inferiority testing},pdfauthor={Minh-Tri Ta, Hoang-Khang Phan, Le-Giang Tran}}
\AtBeginDocument{\justifying}
\fancyhead[L]{\small Cross-fitted forearm electrode reduction}
\fancyhead[C]{}\fancyhead[R]{\small Ta \textit{et al}}


\newcommand{\GRABMyoParticipants}{43}
\newcommand{\GRABMyoSessions}{3}
\newcommand{\GRABMyoGestures}{17}
\newcommand{\GRABMyoActiveGestures}{16}
\newcommand{\GRABMyoWindows}{598{,}689}
\newcommand{\GRABMyoRecords}{15{,}351}
\newcommand{\GRABMyoRecordsFailed}{0}
\newcommand{\GRABMyoAgeMean}{26.4}
\newcommand{\GRABMyoAgeSD}{2.9}
\newcommand{\GRABMyoMale}{23}
\newcommand{\GRABMyoFemale}{20}
\newcommand{\GRABMyoCircMean}{24.1}
\newcommand{\GRABMyoCircSD}{2.3}
\newcommand{\GRABMyoArcSpacing}{3.0}
\newcommand{\GRABMyoActiveWindows}{563{,}472}
\newcommand{\FoldTrainingWindows}{550{,}368}
\newcommand{\TestWindowsPerParticipant}{13{,}104}
\newcommand{\DBFiveSubjects}{10}
\newcommand{\DBFiveAgeMean}{28.0}
\newcommand{\DBFiveAgeSD}{4.0}
\newcommand{\DBFiveMale}{8}
\newcommand{\DBFiveFemale}{2}
\newcommand{\DBFiveWindows}{21{,}089}
\newcommand{\GRABMyoFilesVerified}{30{,}715}
\newcommand{\FullRF}{0.502}
\newcommand{\FullRFLCB}{0.475}
\newcommand{\FullRFUCB}{0.528}
\newcommand{\FullRFSD}{0.090}
\newcommand{\FullRFMin}{0.346}
\newcommand{\FullRFMax}{0.670}
\newcommand{\FullSVM}{0.538}
\newcommand{\FullSVMLCB}{0.513}
\newcommand{\FullSVMUCB}{0.562}
\newcommand{\FullSVMSD}{0.085}
\newcommand{\FullSVMMin}{0.383}
\newcommand{\FullSVMMax}{0.718}
\newcommand{\FullXGB}{0.542}
\newcommand{\FullXGBLCB}{0.517}
\newcommand{\FullXGBUCB}{0.567}
\newcommand{\FullXGBSD}{0.085}
\newcommand{\FullXGBMin}{0.385}
\newcommand{\FullXGBMax}{0.702}
\newcommand{\KStarMarginPtOne}{14}
\newcommand{\KStarMarginPtTwo}{12}
\newcommand{\KStarMarginPtThree}{11}
\newcommand{\KStar}{12}
\newcommand{\KStarReductionPct}{25}
\newcommand{\SVMvsRFDiff}{+0.0362}
\newcommand{\SVMvsRFLCB}{+0.0284}
\newcommand{\SVMvsRFUCB}{+0.0444}
\newcommand{\SVMvsRFW}{15.0}
\newcommand{\SVMvsRFP}{< 0.001}
\newcommand{\SVMvsRFR}{+0.97}
\newcommand{\XGBvsRFDiff}{+0.0408}
\newcommand{\XGBvsRFLCB}{+0.0340}
\newcommand{\XGBvsRFUCB}{+0.0475}
\newcommand{\XGBvsRFW}{10.0}
\newcommand{\XGBvsRFP}{< 0.001}
\newcommand{\XGBvsRFR}{+0.98}
\newcommand{\XGBvsSVMDiff}{+0.0045}
\newcommand{\XGBvsSVMLCB}{-0.0015}
\newcommand{\XGBvsSVMUCB}{+0.0107}
\newcommand{\XGBvsSVMW}{354.0}
\newcommand{\XGBvsSVMP}{= 0.154}
\newcommand{\XGBvsSVMR}{+0.25}
\newcommand{\SVMCalibTwentyK}{0.555}
\newcommand{\SVMCalibFortyK}{0.591}
\newcommand{\CurveKSixteen}{0.541}
\newcommand{\CurveKSixteenSD}{0.084}
\newcommand{\CurveKFourteen}{0.535}
\newcommand{\CurveKFourteenSD}{0.083}
\newcommand{\DeltaKFourteen}{-0.0060}
\newcommand{\DeltaKFourteenLCB}{-0.0081}
\newcommand{\DeltaKFourteenUCB}{-0.0039}
\newcommand{\CurveKThirteen}{0.530}
\newcommand{\CurveKThirteenSD}{0.083}
\newcommand{\DeltaKThirteen}{-0.0109}
\newcommand{\DeltaKThirteenLCB}{-0.0133}
\newcommand{\DeltaKThirteenUCB}{-0.0085}
\newcommand{\CurveKTwelve}{0.528}
\newcommand{\CurveKTwelveSD}{0.084}
\newcommand{\DeltaKTwelve}{-0.0133}
\newcommand{\DeltaKTwelveLCB}{-0.0160}
\newcommand{\DeltaKTwelveUCB}{-0.0106}
\newcommand{\CurveKEleven}{0.521}
\newcommand{\CurveKElevenSD}{0.084}
\newcommand{\DeltaKEleven}{-0.0204}
\newcommand{\DeltaKElevenLCB}{-0.0246}
\newcommand{\DeltaKElevenUCB}{-0.0162}
\newcommand{\CurveKTen}{0.515}
\newcommand{\CurveKTenSD}{0.083}
\newcommand{\DeltaKTen}{-0.0259}
\newcommand{\DeltaKTenLCB}{-0.0311}
\newcommand{\DeltaKTenUCB}{-0.0205}
\newcommand{\CurveKEight}{0.505}
\newcommand{\CurveKEightSD}{0.085}
\newcommand{\DeltaKEight}{-0.0366}
\newcommand{\DeltaKEightLCB}{-0.0440}
\newcommand{\DeltaKEightUCB}{-0.0289}
\newcommand{\CurveKSix}{0.480}
\newcommand{\CurveKSixSD}{0.084}
\newcommand{\DeltaKSix}{-0.0610}
\newcommand{\DeltaKSixLCB}{-0.0713}
\newcommand{\DeltaKSixUCB}{-0.0507}
\newcommand{\CurveKFour}{0.397}
\newcommand{\CurveKFourSD}{0.071}
\newcommand{\DeltaKFour}{-0.1444}
\newcommand{\DeltaKFourLCB}{-0.1590}
\newcommand{\DeltaKFourUCB}{-0.1306}
\newcommand{\CurveKTwo}{0.299}
\newcommand{\CurveKTwoSD}{0.057}
\newcommand{\DeltaKTwo}{-0.2420}
\newcommand{\DeltaKTwoLCB}{-0.2603}
\newcommand{\DeltaKTwoUCB}{-0.2230}
\newcommand{\CurveKOne}{0.181}
\newcommand{\CurveKOneSD}{0.040}
\newcommand{\DeltaKOne}{-0.3603}
\newcommand{\DeltaKOneLCB}{-0.3832}
\newcommand{\DeltaKOneUCB}{-0.3369}
\newcommand{\CurveTrainWindows}{200{,}000}
\newcommand{\LambdaCurveMaxSpread}{0.0040}
\newcommand{\TopFourSet}{4, 6, 7, 14}
\newcommand{\TopFourFolds}{36}
\newcommand{\DroppedFourSet}{9, 11, 12, 16}
\newcommand{\DroppedFourFolds}{43}
\newcommand{\TopFourShare}{34.0}
\newcommand{\BottomEightShare}{39.6}
\newcommand{\UOverAMedian}{2.4}
\newcommand{\UOverAMax}{9.4}
\newcommand{\TopTwelveDiffersLambdaPtTwoFive}{0}
\newcommand{\TopTwelveDiffersLambdaPtFive}{0}
\newcommand{\TopTwelveDiffersLambdaOne}{1}
\newcommand{\HTwoKFourDiff}{+0.0172}
\newcommand{\HTwoKFourLCB}{+0.0071}
\newcommand{\HTwoKFourUCB}{+0.0274}
\newcommand{\HTwoKFourP}{= 0.007}
\newcommand{\HTwoKFourR}{+0.50}
\newcommand{\HTwoKSixDiff}{+0.0370}
\newcommand{\HTwoKSixLCB}{+0.0286}
\newcommand{\HTwoKSixUCB}{+0.0452}
\newcommand{\HTwoKSixP}{< 0.001}
\newcommand{\HTwoKSixR}{+0.92}
\newcommand{\HTwoKEightDiff}{+0.0243}
\newcommand{\HTwoKEightLCB}{+0.0176}
\newcommand{\HTwoKEightUCB}{+0.0308}
\newcommand{\HTwoKEightP}{< 0.001}
\newcommand{\HTwoKEightR}{+0.86}
\newcommand{\HTwoKTwelveDiff}{+0.0074}
\newcommand{\HTwoKTwelveLCB}{+0.0046}
\newcommand{\HTwoKTwelveUCB}{+0.0103}
\newcommand{\HTwoKTwelveP}{< 0.001}
\newcommand{\HTwoKTwelveR}{+0.71}
\newcommand{\HTwoAllSig}{yes}
\newcommand{\PctKFour}{65}
\newcommand{\PctKAboveNinetyFiveFour}{3}
\newcommand{\PctKSix}{89}
\newcommand{\PctKAboveNinetyFiveSix}{11}
\newcommand{\PctKEight}{82}
\newcommand{\PctKAboveNinetyFiveEight}{6}
\newcommand{\PctKTwelve}{66}
\newcommand{\PctKAboveNinetyFiveTwelve}{1}
\newcommand{\VSPermKFourDiff}{-0.0018}
\newcommand{\VSPermKFourP}{= 0.678}
\newcommand{\VSGeomKFourDiff}{-0.0005}
\newcommand{\VSGeomKFourP}{= 0.986}
\newcommand{\VSMiKFourDiff}{+0.1031}
\newcommand{\VSMiKFourP}{< 0.001}
\newcommand{\VSPermKSixDiff}{+0.0126}
\newcommand{\VSPermKSixP}{= 0.009}
\newcommand{\VSGeomKSixDiff}{+0.0341}
\newcommand{\VSGeomKSixP}{< 0.001}
\newcommand{\VSMiKSixDiff}{+0.0934}
\newcommand{\VSMiKSixP}{< 0.001}
\newcommand{\VSPermKEightDiff}{+0.0033}
\newcommand{\VSPermKEightP}{= 0.069}
\newcommand{\VSGeomKEightDiff}{+0.0047}
\newcommand{\VSGeomKEightP}{= 0.267}
\newcommand{\VSMiKEightDiff}{+0.0737}
\newcommand{\VSMiKEightP}{< 0.001}
\newcommand{\VSPermKTwelveDiff}{+0.0001}
\newcommand{\VSPermKTwelveP}{= 0.995}
\newcommand{\VSGeomKTwelveDiff}{+0.0070}
\newcommand{\VSGeomKTwelveP}{= 0.058}
\newcommand{\VSMiKTwelveDiff}{+0.0051}
\newcommand{\VSMiKTwelveP}{= 0.007}
\newcommand{\HThreeMaxAbsDiff}{0.0015}
\newcommand{\HThreeNSig}{0}
\newcommand{\HThreeNTests}{12}
\newcommand{\HThreeMinP}{= 0.216}
\newcommand{\HThreeMinTied}{36}
\newcommand{\FaithProtocolCells}{172}
\newcommand{\FaithProtocolMeanDiff}{0.0014}
\newcommand{\FaithProtocolMaxDiff}{0.0084}
\newcommand{\FaithTopKTwo}{0.263}
\newcommand{\FaithBottomKTwo}{0.292}
\newcommand{\FaithTopKFour}{0.351}
\newcommand{\FaithRandomKFour}{0.366}
\newcommand{\FaithBottomKFour}{0.382}
\newcommand{\FaithTopKSix}{0.376}
\newcommand{\FaithRandomKSix}{0.416}
\newcommand{\FaithBottomKSix}{0.454}
\newcommand{\FaithTopKEight}{0.408}
\newcommand{\FaithRandomKEight}{0.445}
\newcommand{\FaithBottomKEight}{0.469}
\newcommand{\FaithTopKTen}{0.445}
\newcommand{\FaithBottomKTen}{0.475}
\newcommand{\FaithTopKTwelve}{0.455}
\newcommand{\FaithRandomKTwelve}{0.478}
\newcommand{\FaithBottomKTwelve}{0.485}
\newcommand{\FaithTopKFourteen}{0.487}
\newcommand{\FaithBottomKFourteen}{0.491}
\newcommand{\FaithTBKTwoDiff}{-0.0297}
\newcommand{\FaithTBKTwoLCB}{-0.0419}
\newcommand{\FaithTBKTwoUCB}{-0.0166}
\newcommand{\FaithTBKTwoP}{< 0.001}
\newcommand{\FaithTBKFourDiff}{-0.0314}
\newcommand{\FaithTBKFourLCB}{-0.0447}
\newcommand{\FaithTBKFourUCB}{-0.0177}
\newcommand{\FaithTBKFourP}{< 0.001}
\newcommand{\FaithTRKFourDiff}{-0.0151}
\newcommand{\FaithTRKFourLCB}{-0.0239}
\newcommand{\FaithTRKFourUCB}{-0.0061}
\newcommand{\FaithTRKFourP}{= 0.010}
\newcommand{\FaithRBKFourDiff}{-0.0163}
\newcommand{\FaithRBKFourLCB}{-0.0265}
\newcommand{\FaithRBKFourUCB}{-0.0061}
\newcommand{\FaithRBKFourP}{= 0.013}
\newcommand{\FaithTBKSixDiff}{-0.0775}
\newcommand{\FaithTBKSixLCB}{-0.0921}
\newcommand{\FaithTBKSixUCB}{-0.0627}
\newcommand{\FaithTBKSixP}{< 0.001}
\newcommand{\FaithTRKSixDiff}{-0.0398}
\newcommand{\FaithTRKSixLCB}{-0.0487}
\newcommand{\FaithTRKSixUCB}{-0.0307}
\newcommand{\FaithTRKSixP}{< 0.001}
\newcommand{\FaithRBKSixDiff}{-0.0377}
\newcommand{\FaithRBKSixLCB}{-0.0459}
\newcommand{\FaithRBKSixUCB}{-0.0293}
\newcommand{\FaithRBKSixP}{< 0.001}
\newcommand{\FaithTBKEightDiff}{-0.0618}
\newcommand{\FaithTBKEightLCB}{-0.0772}
\newcommand{\FaithTBKEightUCB}{-0.0461}
\newcommand{\FaithTBKEightP}{< 0.001}
\newcommand{\FaithTRKEightDiff}{-0.0377}
\newcommand{\FaithTRKEightLCB}{-0.0476}
\newcommand{\FaithTRKEightUCB}{-0.0273}
\newcommand{\FaithTRKEightP}{< 0.001}
\newcommand{\FaithRBKEightDiff}{-0.0242}
\newcommand{\FaithRBKEightLCB}{-0.0309}
\newcommand{\FaithRBKEightUCB}{-0.0172}
\newcommand{\FaithRBKEightP}{< 0.001}
\newcommand{\FaithTBKTenDiff}{-0.0297}
\newcommand{\FaithTBKTenLCB}{-0.0399}
\newcommand{\FaithTBKTenUCB}{-0.0188}
\newcommand{\FaithTBKTenP}{< 0.001}
\newcommand{\FaithTBKTwelveDiff}{-0.0302}
\newcommand{\FaithTBKTwelveLCB}{-0.0385}
\newcommand{\FaithTBKTwelveUCB}{-0.0218}
\newcommand{\FaithTBKTwelveP}{< 0.001}
\newcommand{\FaithTRKTwelveDiff}{-0.0228}
\newcommand{\FaithTRKTwelveLCB}{-0.0296}
\newcommand{\FaithTRKTwelveUCB}{-0.0158}
\newcommand{\FaithTRKTwelveP}{< 0.001}
\newcommand{\FaithRBKTwelveDiff}{-0.0074}
\newcommand{\FaithRBKTwelveLCB}{-0.0101}
\newcommand{\FaithRBKTwelveUCB}{-0.0046}
\newcommand{\FaithRBKTwelveP}{< 0.001}
\newcommand{\FaithTBKFourteenDiff}{-0.0038}
\newcommand{\FaithTBKFourteenLCB}{-0.0084}
\newcommand{\FaithTBKFourteenUCB}{+0.0009}
\newcommand{\FaithTBKFourteenP}{= 0.150}
\newcommand{\FaithNTests}{15}
\newcommand{\FaithNSig}{14}
\newcommand{\FaithTBAllSig}{no}
\newcommand{\StabSubjPairs}{903}
\newcommand{\StabSubjRho}{0.927}
\newcommand{\StabSubjRhoSD}{0.038}
\newcommand{\StabSubjJacFour}{0.752}
\newcommand{\StabSubjJacSix}{0.905}
\newcommand{\StabSubjJacEight}{0.803}
\newcommand{\StabInnerPairs}{430}
\newcommand{\StabInnerRho}{0.948}
\newcommand{\StabInnerRhoSD}{0.021}
\newcommand{\StabInnerJacFour}{0.816}
\newcommand{\StabInnerJacSix}{0.992}
\newcommand{\StabInnerJacEight}{0.838}
\newcommand{\StabOuterPairs}{903}
\newcommand{\StabOuterRho}{0.987}
\newcommand{\StabOuterRhoSD}{0.008}
\newcommand{\StabOuterJacFour}{0.888}
\newcommand{\StabOuterJacSix}{1.000}
\newcommand{\StabOuterJacEight}{0.923}
\newcommand{\CSOneSOne}{0.492}
\newcommand{\CSOneSOneSD}{0.101}
\newcommand{\CSOneSOneTrain}{183{,}456}
\newcommand{\CSOneSTwo}{0.528}
\newcommand{\CSOneSTwoSD}{0.088}
\newcommand{\CSOneSTwoTrain}{183{,}456}
\newcommand{\CSOneSThree}{0.510}
\newcommand{\CSOneSThreeSD}{0.090}
\newcommand{\CSOneSThreeTrain}{183{,}456}
\newcommand{\CSTwoSThree}{0.530}
\newcommand{\CSTwoSThreeSD}{0.095}
\newcommand{\CSTwoSThreeTrain}{366{,}912}
\newcommand{\CSSTwoVsSOneDiff}{+0.0355}
\newcommand{\CSSTwoVsSOneLCB}{+0.0091}
\newcommand{\CSSTwoVsSOneUCB}{+0.0621}
\newcommand{\CSSTwoVsSOneW}{287.0}
\newcommand{\CSSTwoVsSOneP}{= 0.024}
\newcommand{\CSSThreeVsSOneDiff}{+0.0172}
\newcommand{\CSSThreeVsSOneLCB}{-0.0038}
\newcommand{\CSSThreeVsSOneUCB}{+0.0391}
\newcommand{\CSSThreeVsSOneW}{381.0}
\newcommand{\CSSThreeVsSOneP}{= 0.272}
\newcommand{\CSTwoVsOneDiff}{+0.0210}
\newcommand{\CSTwoVsOneLCB}{+0.0146}
\newcommand{\CSTwoVsOneUCB}{+0.0274}
\newcommand{\CSTwoVsOneW}{77.0}
\newcommand{\CSTwoVsOneP}{< 0.001}
\newcommand{\DBFiveFull}{0.281}
\newcommand{\DBFiveFullSD}{0.099}
\newcommand{\DBFiveFullMin}{0.096}
\newcommand{\DBFiveFullMax}{0.386}
\newcommand{\DBFiveKSixteen}{0.281}
\newcommand{\DBFiveKSixteenSD}{0.099}
\newcommand{\DBFiveKTwelve}{0.268}
\newcommand{\DBFiveKTwelveSD}{0.092}
\newcommand{\DBFiveKEight}{0.252}
\newcommand{\DBFiveKEightSD}{0.068}
\newcommand{\DBFiveKFour}{0.180}
\newcommand{\DBFiveKFourSD}{0.026}
\newcommand{\DBFiveDeltaKEleven}{-0.0211}
\newcommand{\DBFiveDeltaKElevenLCB}{-0.03412}
\newcommand{\DBFiveDeltaKTwelve}{-0.0126}
\newcommand{\DBFiveDeltaKTwelveLCB}{-0.01999}
\newcommand{\DBFiveDeltaKThirteen}{-0.0143}
\newcommand{\DBFiveDeltaKThirteenLCB}{-0.03119}
\newcommand{\DBFiveDeltaKFourteen}{-0.0119}
\newcommand{\DBFiveDeltaKFourteenLCB}{-0.02744}
\newcommand{\DBFiveDeltaKFifteen}{-0.0108}
\newcommand{\DBFiveDeltaKFifteenLCB}{-0.02377}
\newcommand{\DBFiveKStarMarginPtOne}{none}
\newcommand{\DBFiveKStarMarginPtTwo}{12}
\newcommand{\DBFiveKStarMarginPtThree}{12}
\newcommand{\DBFiveNIPassing}{12}
\newcommand{\DBFiveMonteCarloFrac}{95.5}
\newcommand{\DBFiveMonteCarloN}{200}
\newcommand{\DBFiveMonteCarloHits}{191}
\newcommand{\DBFiveMonteCarloLCBMin}{-0.0201}
\newcommand{\DBFiveMonteCarloLCBMax}{-0.0195}
\newcommand{\BipFull}{0.478}
\newcommand{\BipFullSD}{0.089}
\newcommand{\BipKOne}{0.174}
\newcommand{\BipKOneSD}{0.044}
\newcommand{\BipKTwo}{0.240}
\newcommand{\BipKTwoSD}{0.049}
\newcommand{\BipKThree}{0.310}
\newcommand{\BipKThreeSD}{0.070}
\newcommand{\BipKFour}{0.356}
\newcommand{\BipKFourSD}{0.080}
\newcommand{\BipKFive}{0.405}
\newcommand{\BipKFiveSD}{0.085}
\newcommand{\BipKSix}{0.442}
\newcommand{\BipKSixSD}{0.086}
\newcommand{\BipKSeven}{0.464}
\newcommand{\BipKSevenSD}{0.085}
\newcommand{\BipKEight}{0.480}
\newcommand{\BipKEightSD}{0.089}
\newcommand{\BipDeltaKOne}{-0.3052}
\newcommand{\BipDeltaKOneLCB}{-0.3298}
\newcommand{\BipDeltaKTwo}{-0.2399}
\newcommand{\BipDeltaKTwoLCB}{-0.2615}
\newcommand{\BipDeltaKThree}{-0.1697}
\newcommand{\BipDeltaKThreeLCB}{-0.1852}
\newcommand{\BipDeltaKFour}{-0.1240}
\newcommand{\BipDeltaKFourLCB}{-0.1392}
\newcommand{\BipDeltaKFive}{-0.0742}
\newcommand{\BipDeltaKFiveLCB}{-0.0845}
\newcommand{\BipDeltaKSix}{-0.0375}
\newcommand{\BipDeltaKSixLCB}{-0.0451}
\newcommand{\BipDeltaKSeven}{-0.0158}
\newcommand{\BipDeltaKSevenLCB}{-0.0203}
\newcommand{\BipKStarMarginPtOne}{none}
\newcommand{\BipKStarMarginPtTwo}{none}
\newcommand{\BipKStarMarginPtThree}{7}
\newcommand{\BipolarMonteCarloFracSeven}{0}
\newcommand{\BipolarMonteCarloHitsSeven}{0}
\newcommand{\BipolarMonteCarloN}{200}
\newcommand{\BipolarMonteCarloLCBMaxSeven}{-0.0203}
\newcommand{\BipolarMonteCarloLCBMinSeven}{-0.0207}
\newcommand{\BipCurveRefK}{0.480}
\newcommand{\BipTopThreeSet}{5, 6, 7}
\newcommand{\BipTopThreeFolds}{30}
\newcommand{\FixedDev}{22}
\newcommand{\FixedTest}{21}
\newcommand{\FixedChannels}{1, 2, 4, 5, 6, 7, 8, 10, 13, 14, 15, 16}
\newcommand{\FixedDropped}{3, 9, 11, 12}
\newcommand{\FixedFull}{0.522}
\newcommand{\FixedTwelve}{0.510}
\newcommand{\FixedDiff}{-0.0124}
\newcommand{\FixedLCB}{-0.0179}
\newcommand{\FixedUCB}{-0.0077}
\newcommand{\FixedLamZeroSame}{yes}
\newcommand{\FixedNLambdaSameOrig}{4}
\newcommand{\FixedSwapDev}{21}
\newcommand{\FixedSwapTest}{22}
\newcommand{\FixedSwapDropped}{9, 11, 12, 13}
\newcommand{\FixedSwapNLambdaSame}{4}
\newcommand{\FixedSwapFull}{0.529}
\newcommand{\FixedSwapTwelve}{0.516}
\newcommand{\FixedSwapDiff}{-0.0135}
\newcommand{\FixedSwapLCBSix}{-0.020020}
\newcommand{\FixedSwapUCB}{-0.0066}
\newcommand{\FixedSwapNI}{no}
\newcommand{\FixedDroppedBoth}{9, 11, 12}
\newcommand{\FixedDroppedBothN}{3}
\newcommand{\RedundDiff}{-0.0005}
\newcommand{\RedundLCB}{-0.0021}
\newcommand{\RedundUCB}{+0.0010}
\newcommand{\RedundRho}{0.994}
\newcommand{\RedundTopTwelveSame}{34}
\newcommand{\KStarHolmPtOne}{14}
\newcommand{\KStarFixSeqPtOne}{14}
\newcommand{\KStarHolm}{12}
\newcommand{\KStarFixSeq}{12}
\newcommand{\KStarHolmPtThree}{11}
\newcommand{\KStarFixSeqPtThree}{11}
\newcommand{\HolmPKTwelve}{= 0.002}
\newcommand{\HolmPKEleven}{= 1.000}
\newcommand{\HolmSameAllLambda}{yes}
\newcommand{\DBFiveKStarHolmPtOne}{none}
\newcommand{\DBFiveKStarHolmPtTwo}{none}
\newcommand{\DBFiveKStarHolmPtThree}{12}
\newcommand{\DBFiveDeltaKTwelveLCBSix}{-0.019994}
\newcommand{\MarginPctOfRef}{3.7}
\newcommand{\MarginPctAboveChance}{4.2}
\newcommand{\LOOMonoFolds}{43}
\newcommand{\LOOMonoIdentical}{43}
\newcommand{\LOOMonoInProbe}{6}
\newcommand{\LOODBFiveFolds}{10}
\newcommand{\LOODBFiveIdentical}{10}
\newcommand{\LOODBFiveInProbe}{10}
\newcommand{\LOOBipFolds}{43}
\newcommand{\LOOBipIdentical}{43}
\newcommand{\LOOBipInProbe}{6}
\newcommand{\LOOMonoMin}{8.52}
\newcommand{\LOOMonoMax}{9.45}
\newcommand{\LOOBipPooled}{3.40}
\newcommand{\LOOBipRule}{4}
\newcommand{\HTwoSameKFour}{40}
\newcommand{\HTwoSameKSix}{43}
\newcommand{\HTwoSameKEight}{40}
\newcommand{\HTwoSameKTwelve}{43}

\newcommand{\RevDBFiveBandLow}{10}
\newcommand{\RevDBFiveBandHigh}{89}
\newcommand{\RevDBFiveBandRequestedHigh}{95}
\newcommand{\RevDBFiveFeatureRows}{21089}
\newcommand{\RevShapDelta}{-0.0133}
\newcommand{\RevShapLow}{-0.0160}
\newcommand{\RevShapHigh}{-0.0106}
\newcommand{\RevShapMean}{0.528}
\newcommand{\RevGeomDelta}{-0.0193}
\newcommand{\RevGeomLow}{-0.0254}
\newcommand{\RevGeomHigh}{-0.0132}
\newcommand{\RevGeomMean}{0.522}
\newcommand{\RevPrimaryFull}{0.541}
\newcommand{\RevGeomContrast}{+0.0060}
\newcommand{\RevGeomContrastLow}{-0.0010}
\newcommand{\RevGeomContrastHigh}{+0.0127}
\newcommand{\RevGeomContrastP}{0.024}
\newcommand{\RevLossCount}{11}
\newcommand{\RevImprovedCount}{1}
\newcommand{\RevWorstParticipant}{27}
\newcommand{\RevMedianLoss}{-0.0125}
\newcommand{\RevWorstLoss}{-0.0361}
\newcommand{\RevBestChange}{+0.0144}
\newcommand{\RevLossSD}{+0.0093}
\newcommand{\RevGestureLossCount}{4}
\newcommand{\RevWorstGestureDelta}{-0.0456}
\newcommand{\RevWorstGestureLow}{-0.0594}
\newcommand{\RevWorstGestureHigh}{-0.0323}
\newcommand{\RevMatchedSixDiff}{+0.0309}
\newcommand{\RevMatchedSixLow}{+0.0211}
\newcommand{\RevMatchedSixHigh}{+0.0410}
\newcommand{\RevMatchedSixP}{< 0.001}
\newcommand{\RevMatchedTwelveDiff}{+0.0002}
\newcommand{\RevMatchedTwelveLow}{-0.0029}
\newcommand{\RevMatchedTwelveHigh}{+0.0034}
\newcommand{\RevMatchedTwelveP}{= 0.976}
\newcommand{\RevPenaltyBenefitCount}{0}
\newcommand{\RevPenaltyHarmCount}{0}
\newcommand{\RevCostTraining}{117.1}
\newcommand{\RevCostShap}{4.0}
\newcommand{\RevCostPermutation}{10.6}
\newcommand{\FixedLayoutMonteCarloHits}{90}
\newcommand{\FixedLayoutMonteCarloN}{200}
\newcommand{\FixedLayoutMonteCarloPercent}{45}
\newcommand{\RevGeomKStarCI}{15}
\newcommand{\RevGeomKStarNull}{15}

\renewcommand{\topfraction}{0.9}\renewcommand{\bottomfraction}{0.85}\renewcommand{\textfraction}{0.08}\renewcommand{\floatpagefraction}{0.75}
\widowpenalty=10000\clubpenalty=10000\displaywidowpenalty=10000
\begin{document}
\title{Subject-independent forearm sEMG electrode reduction with cross-fitted attribution and non-inferiority testing}
\author{Minh-Tri Ta$^1$, Hoang-Khang Phan$^{2,3}$ and Le-Giang Tran$^{4,*}$}

\affil{$^1$College of Engineering and Computer Science, VinUniversity, Hanoi, Vietnam}

\affil{$^2$Department of Biomedical Engineering, Faculty of Applied Science, Ho Chi Minh City University of Technology (HCMUT), 268 Ly Thuong Kiet Street, District 10, Ho Chi Minh City, Vietnam}

\affil{$^3$Vietnam National University Ho Chi Minh City, Linh Trung Ward, Thu Duc City, Ho Chi Minh City, Vietnam}

\affil{$^4$Faculty of Electronics Technology, Industrial University of Ho Chi Minh City, Ho Chi Minh City, Vietnam}

\affil{$^*$Author to whom any correspondence should be addressed.}

\email{tranlegiang@iuh.edu.vn}

\keywords{surface electromyography, electrode selection, cross-fitting, SHAP, non-inferiority, participant-independent validation}
\begin{abstract}
\textit{Objective.} Determine how many forearm surface electromyography electrodes can be removed while preserving mean gesture recognition for unseen participants, and assess the additional value of attribution-based selection.
\textit{Approach.} We analyzed \GRABMyoParticipants\ GRABMyo participants using sixteen monopolar electrodes and sixteen active gestures. Participant-disjoint inner cross-fitting estimated TreeSHAP rankings within each outer leave-one-subject-out fold. Reduced models were retrained and compared with full16 using participant-bootstrap intervals and a predefined 0.02 Macro-F1 engineering tolerance. Exploratory analyses included primary-budget geometric selection, paired gesture predictions, shared-resource permutation selection, and task, window and classifier sensitivities.
\textit{Main results.} SHAP12 achieved mean Macro-F1 \RevShapMean\ versus \RevPrimaryFull\ for full16; its change was \RevShapDelta\ (95\% CI \RevShapLow\ to \RevShapHigh). The smallest qualifying count was \KStarMarginPtTwo. Nevertheless, \RevLossCount\ participants lost more than 0.02 and \RevGestureLossCount\ gestures showed corrected declines. Geometric12 did not establish non-inferiority, and the direct SHAP12 mean advantage remained uncertain (\RevGeomContrast; 95\% CI \RevGeomContrastLow\ to \RevGeomContrastHigh). Attribution outperformed mean random subsets with the reduced-budget classifier. Shared-resource SHAP outperformed grouped permutation at six electrodes, but its twelve-electrode advantage remained uncertain. Stability regularization did not establish an improvement.
\textit{Significance.} One quarter of this monopolar array could be removed within the mean tolerance under offline participant-independent evaluation. This validates a selection procedure rather than one locked physical layout; individual losses and uncertain transfer require independent validation before device claims.
\end{abstract}

\section{Introduction}
Surface electromyography (sEMG) provides a non-invasive signal for gesture recognition and myoelectric control, but the recording interface must balance spatial information against the burden of additional electrodes \citep{scheme2011electromyogram,farina2014extraction,vujaklija2016new}. Fewer recording sites could simplify an interface. Whether that reduction preserves recognition for an unseen participant is an empirical question, and offline recognition alone does not establish a usable control interface.

Channel selection has been studied with feature filters, wrappers and spatial criteria \citep{phinyomark2012feature,stango2015spatial,qu2021reduce}. SHAP-guided EMG channel and feature optimization is also established \citep{rani2025shap}; related work uses Shapley explanations for similar sEMG gestures \citep{wang2024shap} and explainable feature selection with fusion \citep{gehlot2024fusion}. These contributions make attribution-based reduction a relevant comparator rather than a new algorithmic premise. Their performance estimates cannot be directly compared with this study across different tasks, sensors and validation protocols.

For a system intended to work across people, selection must exclude the eventual test participant. Using that participant to standardize features, estimate importance or choose a subset can produce optimistic estimates. Participant-disjoint inner evaluation also avoids measuring attribution solely on observations used to fit the attributing model. This motivates an outer leave-one-subject-out (LOSO) evaluation in which each entire selection procedure is reconstructed from the outer training participants.

Attribution magnitude describes predictive reliance of a fitted model, not a direct measurement of muscle activation. TreeSHAP distributes a model's prediction over its input features \citep{lundberg2017unified,lundberg2020local}; aggregation over an electrode's features offers a ranking, but does not turn it into a coalition-level Shapley value for that electrode. Correlated features and variation across participants can affect this ranking. Attribution stability has already motivated resampling and regularization methods \citep{visani2022statistical,xiang2026roshap}; a stability penalty therefore requires validation rather than a novelty claim.

Electrode count requires an explicit tolerance for performance loss. Non-inferiority asks whether a reduced system remains within such a tolerance, whereas a non-significant superiority test does not establish preservation \citep{schuirmann1987comparison,walker2011understanding}. A mean tolerance is also different from preserving every participant or gesture. Reporting only an aggregate curve could hide losses that matter to an intended application.

We evaluate cross-fitted electrode reduction on GRABMyo and separately examine NinaPro DB5. The primary question is the minimum count supported by a specified engineering margin. Secondary questions concern utility beyond random and spatial selection, stability regularization, and deletion with retraining. Exploratory analyses use a geometric subset under the primary training protocol, paired participant and gesture predictions, and shared-resource SHAP and grouped permutation rankings. Sensitivities and split-derived layouts delimit the claim: the LOSO experiment evaluates a selection procedure, not a prospectively validated common physical layout.

\begin{figure}[htbp]\centering
\includegraphics[width=\textwidth]{../figures/vector/fig1_montage.pdf}
\caption{Recording topology and task. The source montage places F1--F8 in the proximal ring and F9--F16 in the distal ring, separated by 2 cm. Corresponding circumferential positions are aligned. The forearm outline is schematic; recording sites were not registered to individual muscles. The unwrapped map preserves electrode numbers without assigning an absolute anatomical angle.}
\label{fig:montage}\end{figure}

\section{Methods}
\subsection{Datasets, selection units and signal processing}
GRABMyo v1.1.0 contains \GRABMyoParticipants\ able-bodied participants recorded over \GRABMyoSessions\ days with an EMGUSB2+ amplifier at 2048 Hz \citep{pradhan2022multiday,jiang2024grabmyo}. Source-reported demographics were \GRABMyoMale\ male and \GRABMyoFemale\ female, age \GRABMyoAgeMean\ $\pm$ \GRABMyoAgeSD\ years. The forearm montage has sixteen monopolar sites in two rings (figure~\ref{fig:montage}); each site supplies one channel, so electrode and channel selection coincide here. Counts refer only to these signal recording sites; reference and ground requirements are not evaluated. The primary task uses sixteen active gestures, excluding rest. All \GRABMyoFilesVerified\ distributed files matched published checksums, and \GRABMyoRecords\ records were processed. Source demographics and the official gesture dictionary are provided with the reproducibility package.

Signals were filtered at 20--450 Hz using a Butterworth band-pass with SciPy order argument 4 (eight poles per pass), applied forward and backward for zero phase, followed by a forward--backward 50 Hz notch ($Q=30$). They were then divided into 250 ms windows with 50\% overlap within trial boundaries. The non-causal filter limits the analysis to offline use. Each channel supplied mean absolute value, RMS, variance, waveform length, zero crossings, slope sign changes, integrated EMG, Willison amplitude, mean frequency and median frequency. The active-task cache contains \GRABMyoActiveWindows\ windows with 160 features each.

Count-feature thresholds followed a common rule: 1.15 times the median absolute signal first difference, rounded to one significant figure. The monopolar threshold ($10^{-2}$ mV) was calibrated from 54 unfiltered records from six participants before participant partitioning. Excluding each participant left the rounded threshold unchanged in all \LOOMonoFolds\ folds. Threshold calibration therefore involved pooled development data, although its value was invariant to participant exclusion. Fold-specific feature means and standard deviations, classifier fitting and attribution ranking exclude the outer test participant.

NinaPro DB5 was analyzed separately: \DBFiveSubjects\ participants, two Myo armbands, sixteen channels sampled at 200 Hz and exercise 1's twelve finger movements \citep{pizzolato2017comparison}. Refined labels determine movement segments. Filtering used an effective passband of \RevDBFiveBandLow--\RevDBFiveBandHigh\ Hz without a notch. The 250 ms windows contain 50 samples. The threshold is seven quantization steps in the stored integer representation; every participant-exclusion check reproduced it. Source-reported DB5 demographics were \DBFiveMale\ male and \DBFiveFemale\ female, age \DBFiveAgeMean\ $\pm$ \DBFiveAgeSD\ years. Sex is reported as recorded in the source datasets; no sex-specific effect is estimated. No anatomical correspondence between DB5 and GRABMyo channel numbers is assumed.

\paragraph{Ethics} This is a secondary analysis of publicly distributed, de-identified recordings, with no new recruitment or recording. Original study approvals and consent are described in the dataset publications.

\subsection{Outer evaluation and classifier budgets}
Each outer LOSO fold reserves all observations of one participant $p$. For cross-fitted attribution, an inner five-fold GroupKFold partitions only the other participants (figure~\ref{fig:workflow}). Inner feature standardization uses inner-training observations only; final reduced models use the corresponding outer-training statistics. No outer test observations enter fitting or rankings.

The primary reduction curve uses XGBoost with 300 boosting rounds, maximum depth 6, learning rate 0.1, row and column subsampling 0.8, histogram tree construction. Its multiclass soft-probability objective produces probabilities whose argmax gives the predicted class. For the primary curve, final models across all counts and fixed penalty rankings within an outer fold use the same stratified sample of 200,000 outer-training windows. Models are retrained from scratch on the selected features. Selected features retain their original column order.

Training budgets differ across comparisons: the full-array XGBoost classifier context uses all training windows, the primary curve uses 200,000, and the random-selection comparison uses 50 boosting rounds and 150,000. The full-array SVM uses 40,000 windows; random forest uses 200,000. These are budget-limited comparators, not proven lower bounds or matched classifier-capacity experiments. Budgets were calibrated using two outer-fold outcomes. The analysis design followed pilot work and was not independently preregistered. This development exposure limits prospective interpretation.

\begin{figure}[htbp]\centering
\includegraphics[width=\textwidth]{../figures/vector/fig2_crossfit.pdf}
\caption{Participant access in the nested procedure. Each inner attributing model excludes the participants on whom its attributions are evaluated. The outer participant is reserved until final reduced-model testing. This diagram describes fold-specific scaling, fitting and ranking; pooled threshold calibration and budget calibration are separate development exposures documented in the text.}
\label{fig:workflow}\end{figure}

\subsection{Cross-fitted attribution and its participant-level extension}
Inner XGBoost models use at most 200,000 training windows sampled uniformly without replacement after participant splitting. Their standardization statistics use the full inner-training pool, not the capped model-fit sample. Exact tree-path-dependent TreeSHAP contributions are computed on raw class margins for at most 20,000 inner-held-out windows. For window $i$ and electrode $c$, absolute contributions are averaged across classes $g$ and its ten features $f$:
\begin{equation}\Phi_{i,c}=\frac{1}{GF}\sum_{g=1}^G\sum_{f=1}^F|\phi_{i,g,c,f}|.\end{equation}
The mean over inner-held-out windows is normalized across electrodes to give the inner-fold score $I_{j,c}$. The central magnitude and unscaled median absolute deviation are
\begin{equation}A_c=\operatorname{median}_j I_{j,c},\quad U_c=\operatorname{median}_j|I_{j,c}-A_c|,\quad Q_c(\lambda)=A_c-\lambda U_c.\end{equation}
All fixed penalties $\lambda\in\{0,0.25,0.5,1\}$ are reported, with lower electrode index breaking ties. No penalty is chosen using outer scores. The unpenalized ranking is called cross-fitted SHAP; the positive penalties assess stability regularization.

Pooling people within five inner folds can suppress person-level heterogeneity. A post-hoc extension therefore computes the median and MAD over normalized individual attribution vectors of all 42 outer-training participants. These vectors share five fitted inner models; their dispersion is descriptive heterogeneity, not a standard error. All penalties and counts are evaluated under the primary protocol. A post-hoc descriptive analysis compared selected physical site sets and Jaccard similarity at every fixed penalty and count. It assesses whether penalties change the physical inputs and does not select a penalty.

\subsection{Primary endpoint and participant-level inference}
Macro-F1 is the unweighted mean of the gesture F1 scores for each held-out participant. For count $k$, the paired change is $\Delta_{p,k}=\mathrm{F1}_{p,k}-\mathrm{F1}_{p,16}$. The predefined primary engineering tolerance is $\delta=0.02$ on the 0--1 scale. The primary decision rule requires the lower 95\% percentile-bootstrap bound on the participant mean change to exceed $-\delta$; $k^*$ is the smallest qualifying count below sixteen. Sensitivity margins are 0.01 and 0.03. This criterion concerns mean recognition performance, has no established clinical interpretation and does not imply equivalence or individual preservation.

The participant is the statistical unit: 43 for GRABMyo, ten for DB5 and 21 or 22 test participants for split-layout checks. Bootstrap intervals use 10,000 resamples of paired participant scores. They are conditional on the observed fitted LOSO models. Overlapping training cohorts and the lack of refitting within every bootstrap mean these intervals do not capture all new-cohort training uncertainty \citep{bengio2004variance}.

The percentile-tail Holm procedure uses the fraction of unshifted bootstrap mean differences at or below $-\delta$, with a Monte Carlo floor. This is a percentile-tail quantity rather than a bootstrap distribution generated under the null. An exploratory sensitivity uses a null-centred studentized bootstrap \citep{davison1997bootstrap}. For $d_p=\Delta_{p,k}$, residuals $d_p-\bar d$ are resampled, their studentized means are compared with $(\bar d+\delta)/(s_d/\sqrt n)$, and the upper-tail probability uses a plus-one correction. Holm adjustment covers the fifteen reduced-count tests at one-sided $\alpha=0.025$, separately for each condition, margin and penalty.

Paired secondary contrasts use Wilcoxon signed-rank tests and participant mean-change percentile intervals. These address different estimands: signed ranks are not a test of the mean. Holm families are stated explicitly; a non-significant contrast is not evidence of equivalence. The secondary test families are post-hoc and are not prospectively confirmatory.

\subsection{Selection comparators and attribution diagnostics}
Reduced-budget comparisons use random subsets, cross-fitted SHAP, stability-regularized SHAP ($\lambda=0.5$), feature-wise permutation, mutual information and geometric farthest-point selection. Random selection uses 100 subsets per participant at counts 4, 6 and 8, and 500 at count 12. Eight attribution-versus-mean-random contrasts receive Holm correction. Feature-wise permutation uses one model fitted on 20,000 windows and three feature-wise permutations on 8,000 non-overlapping windows from the same outer-training participant pool. Its 40,000-window ranking pool was standardized before the window split; this was not participant-disjoint inner validation. The outer test participant remained excluded. SHAP instead uses the participant-disjoint cross-fitting above. Thus shared final evaluation does not equalize their selection resources or validation designs.

Exploratory geometric selection uses the primary 300-round, 200,000-window protocol, the same outer training sample and full16 reference. Geometric selection uses the fixed two-ring topology, nominal 24 cm circumference and 2 cm separation, starts at F1 and successively maximizes distance from selected sites; ties favor the lower index. It is not tuned to participant outcomes. The complete geometric count curve evaluates this ordering under the primary protocol. It is a post-hoc analysis motivated by the twelve-electrode comparison.

A separate exploratory resource-matched comparison shares five inner models between SHAP and grouped-electrode permutation, with up to 200,000 uniformly sampled training and 2,000 validation windows per fold. Equal validation quotas represent every inner-held-out participant. All ten features of one electrode are permuted jointly within each participant, three repeats, scoring the equally weighted mean participant Macro-F1 loss. Signed losses are retained, normalized by their summed absolute magnitude, and aggregated by the median across inner folds. SHAP uses the identical observations and normalizes absolute contribution magnitudes. Final models at counts 6 and 12 use the primary protocol; two contrasts receive Holm correction. Measured training and selection times are reported separately. Runtime measurements include concurrent workstation load and are descriptive rather than an isolated benchmark. Because permutation grouping, participant-disjoint validation, scaler isolation and resources all differ from feature-wise permutation, this experiment does not isolate a budget effect.

Deletion with retraining uses the unpenalized primary ordering and compares removal of high-ranked, random and low-ranked electrodes using the lighter model. Its fifteen contrasts form one Holm family. This diagnostic tests predictive ranking relevance; it cannot identify causal muscles or validate an online controller.

\subsection{Participant and gesture degradation}
Primary full16, SHAP12 and geometry12 predictions are indexed to identical outer-test observations. Participant losses, quantiles, worst loss and the number exceeding the primary margin are descriptive. Gesture-specific F1 differences are paired over all participants, with Holm correction across sixteen gestures. Confusion changes average row-normalized confusion matrices within participants so people, rather than windows, receive equal weight. Gesture names follow the explicit numeric ordering in MotionSequence.txt.

\subsection{Sensitivity, transfer and fixed-layout scope}
Rest-inclusive recognition and 150/500 ms windows use raw-signal re-extraction and complete count curves with the primary 250 ms ranking held fixed. The 150 ms window is 307 samples with stride 153; 250/500 ms use 512/1024 samples and half-window strides. The same 200,000-window count does not equate recording duration across windows. A random-forest sensitivity uses a 200-tree, minimum-leaf-size-2 configuration at counts 12 and 16, the unchanged XGBoost ranking and the same training sample. It tests those counts only and cannot establish a random-forest minimum count.

DB5 applies the procedure separately with its own sixteen-channel reference. Derived bipolar signals $B_j=F_j-F_{j+8}$ use eight aligned pairs, re-extracted features and the eight-pair reference; they are not physical bipolar-sensor measurements. Full-array cross-session analyses compare training on session 1 with testing on sessions 1--3, and training on sessions 1--2 with testing on session 3. They do not test a reduced layout under displacement, and the latter comparison confounds data quantity with session diversity.

Fixed-layout checks use odd-ID development/even-ID testing and the reverse, each deriving a layout using development participants only and the five-fold median/MAD aggregation rule. Their inner and final models use all development windows with training-only standardization and the same 300-round classifier, rather than the primary curve's 200,000-window cap. The pooled feature thresholds also apply to these checks. The count of twelve had already been selected using the full cohort, so neither direction independently validates count selection. The layouts differ and must be reported separately. A post-hoc Monte Carlo sensitivity analysis repeated the percentile bootstrap under 200 fixed randomizations, with 10,000 resamples per repetition. The distribution of decisions quantifies numerical sensitivity near the boundary; it does not estimate the probability that non-inferiority is true. No fresh, independently collected cohort with this montage was available. Feature redundancy and ranking-agreement diagnostics are documented in the Supplementary; high outer-fold agreement is qualified by extensive training-set overlap.

\section{Results}
\subsection{How many electrodes meet the mean performance tolerance?}
Under the primary curve protocol, full16 achieved mean Macro-F1 \RevPrimaryFull\ and SHAP12 achieved \RevShapMean. The smallest count meeting the 0.02 margin was \KStarMarginPtTwo, with paired change \RevShapDelta\ (95\% CI \RevShapLow\ to \RevShapHigh). Eleven did not meet this criterion: change \DeltaKEleven\ (\DeltaKElevenLCB\ to \DeltaKElevenUCB). The primary percentile decision and the exploratory null-centred, Holm-adjusted sensitivity gave the same primary minimum count: \KStarMarginPtOne, \KStarMarginPtTwo\ and \KStarMarginPtThree\ at margins 0.01, 0.02 and 0.03, respectively. The four penalty curves did not change these count decisions (figure~\ref{fig:ni}; full results in the Supplementary).

Mean preservation did not extend uniformly to individuals. For SHAP12 minus full16, the median change was \RevMedianLoss\ and participant SD was \RevLossSD. \RevImprovedCount\ participant improved, while \RevLossCount\ lost more than 0.02. The worst observed change was \RevWorstLoss\ (participant \RevWorstParticipant), and the best was \RevBestChange. The margin is a mean-level engineering criterion, so these counts do not redefine individual non-inferiority.

\begin{table}[htbp]
\centering\small
\caption{Matched primary protocol: reduced subset minus full16. Each row uses all 43 paired participants, 300 boosting rounds and 200,000 training windows.}\label{tab:primary}
\begin{adjustbox}{max width=\textwidth}
\begin{tabular}{lllll}\toprule
Subset & Mean F1 & $\bar\Delta$ & 95\% CI & NI at 0.02\\\midrule
SHAP12 & 0.528 & -0.0133 & [-0.0160, -0.0106] & yes\\
Geometry12 & 0.522 & -0.0193 & [-0.0254, -0.0132] & no\\
\bottomrule\end{tabular}
\end{adjustbox}
\end{table}

\begin{figure}[htbp]\centering
\includegraphics[width=\textwidth]{../figures/vector/fig3_noninferiority.pdf}
\caption{Count reduction and its limits. (A) Participant mean curves for all fixed penalties; shading shows the 95\% participant-bootstrap interval for the unpenalized curve. (B) Mean paired change against the same curve's full16 reference, with 95\% intervals and the 0.02 tolerance. (C) All participant changes at SHAP12, ordered by change; orange identifies observed losses larger than 0.02. The horizontal tolerance in C is a descriptive reference, not an individual inference threshold.}
\label{fig:ni}\end{figure}

\subsection{What does attribution add beyond random and spatial selection?}
With the reduced-budget classifier, attribution-ranked subsets outperformed the participant's mean random-subset baseline at all four tested counts after correction. For the stability-regularized ranking, the mean advantages at counts 4, 6, 8 and 12 were \HTwoKFourDiff, \HTwoKSixDiff, \HTwoKEightDiff\ and \HTwoKTwelveDiff. These absolute scores belong to the lighter comparison model and should not be substituted for the primary non-inferiority curve.

With the primary model and training budget matched, geometric12 scored \RevGeomMean, with change \RevGeomDelta\ (95\% CI \RevGeomLow\ to \RevGeomHigh) from full16. Its lower bound crossed the 0.02 tolerance. The direct SHAP12-minus-geometric12 mean contrast was \RevGeomContrast\ (\RevGeomContrastLow\ to \RevGeomContrastHigh), whose interval included zero. The nominal signed-rank $p=\RevGeomContrastP$ addresses ranks rather than the mean; it does not establish a positive mean advantage. The full fixed geometric curve had minimum qualifying counts of \RevGeomKStarCI\ under the percentile rule and \RevGeomKStarNull\ under the null-centred Holm sensitivity at the 0.02 margin; the full curve and decision table are provided in the Supplementary. Different non-inferiority decisions do not establish a difference between methods \citep{gelman2006difference}.

The shared-model comparison is reported separately in table~\ref{tab:matched} and figure~\ref{fig:utility}(C). At counts 6 and 12, SHAP-minus-grouped-permutation mean changes were \RevMatchedSixDiff\ (\RevMatchedSixLow\ to \RevMatchedSixHigh) and \RevMatchedTwelveDiff\ (\RevMatchedTwelveLow\ to \RevMatchedTwelveHigh), with Holm $p\RevMatchedSixP$ and $p\RevMatchedTwelveP$. Mean measured per-outer-fold times were \RevCostTraining\ s for shared inner training, \RevCostShap\ s for SHAP selection and \RevCostPermutation\ s for grouped permutation. These hardware-dependent measurements describe the executed comparison rather than online latency.

At $k=6$, the mean interval supports a positive SHAP advantage; the corrected signed-rank contrast is significant. At $k=12$, the mean advantage remains uncertain because its interval crosses zero; the corrected signed-rank contrast is not statistically detectable.

\begin{table}[htbp]
\centering\small
\caption{Exploratory SHAP minus grouped-electrode permutation. The two methods share the same five inner models and validation observations per outer fold. Holm correction covers the two contrasts.}\label{tab:matched}
\begin{adjustbox}{max width=\textwidth}
\begin{tabular}{llllll}\toprule
$k$ & SHAP & Permutation & $\bar\Delta$ & 95\% CI & $p_{\rm Holm}$\\\midrule
6 & 0.480 & 0.449 & +0.0309 & [+0.0211, +0.0410] & $<0.001$\\
12 & 0.528 & 0.528 & +0.0002 & [-0.0029, +0.0034] & 0.976\\
\bottomrule\end{tabular}
\end{adjustbox}
\end{table}

\begin{figure}[htbp]\centering
\includegraphics[width=\textwidth]{../figures/vector/fig4_selection_utility.pdf}
\caption{Incremental selection utility under distinct protocols. (A) Six-method comparison with a common lighter final classifier (50 boosting rounds, 150,000 training windows); selection-stage resources and validation designs differed, and penalized SHAP uses $\lambda=0.5$. (B) Direct SHAP12-minus-geometric12 mean contrast under the primary protocol (300 boosting rounds, 200,000 windows). (C) SHAP-minus-grouped-permutation contrasts with shared inner models and validation observations and primary-protocol final models. Error bars in B--C are participant mean-change bootstrap intervals. Panels are not a common absolute-score benchmark.}
\label{fig:utility}\end{figure}

\subsection{Which predictive information is lost, and does regularization help?}
Attribution has consistent spatial preferences but does not register sites to individual muscles (figure~\ref{fig:task}(A)). The same top-four set (\TopFourSet) occurred in \TopFourFolds\ outer folds. Outer-fold agreement is descriptive because training cohorts overlap heavily.

The gesture analysis found \RevGestureLossCount\ corrected losses among sixteen active gestures. The largest mean decline was gesture 14, forearm pronation: \RevWorstGestureDelta\ (95\% CI \RevWorstGestureLow\ to \RevWorstGestureHigh). Gestures 2 (thumb adduction), 9 (index finger extension) and 10 (thumb finger extension) also showed corrected declines. These paired gesture-F1 results expose task-specific losses hidden by the global mean; no anatomical cause is inferred from them. Gesture IDs are distinct from electrode indices.

Deletion with retraining produced the expected ordering for the common random-comparison counts: removing highly ranked electrodes gave lower performance than random removal, which gave lower performance than removal of weakly ranked electrodes. At six electrodes retained, Remove-Top minus Remove-Bottom was \FaithTBKSixDiff\ (\FaithTBKSixLCB\ to \FaithTBKSixUCB; Holm $p\FaithTBKSixP$). This supports predictive relevance under retraining, rather than a causal physiological interpretation.

None of the \HThreeNTests\ penalty contrasts showed corrected improvement over magnitude alone. The participant-level extension, which exposes individual heterogeneity instead of pooling it into five groups, yielded \RevPenaltyBenefitCount\ corrected beneficial and \RevPenaltyHarmCount\ corrected harmful penalty contrasts across the twelve fixed comparisons (the Supplementary penalty table). All penalties remain reported; outer-test performance was not used to choose one.

At twelve electrodes, every fixed participant-level penalty selected the same physical site set as its unpenalized counterpart in all 43 outer folds. These identical inputs give identical scores under the fixed classifier protocol; this is not an equivalence test of different layouts.


\begin{figure}[htbp]\centering
\includegraphics[width=\textwidth]{../figures/vector/fig5_task_and_attribution.pdf}
\caption{What the ranking represents and what reduction loses. (A) Mean fold-median normalized absolute attribution over outer folds on the documented topology; the scale is in percent and measures model relevance, not muscle activation. (B) Gesture-specific mean F1 changes with participant-bootstrap intervals; orange marks Holm-adjusted signed-rank $p<0.05$ across sixteen gestures. Names follow MotionSequence.txt, with paired-finger names shortened for display. (C) Deletion with retraining under the lighter model. (D) Participant-level aggregation, penalized minus unpenalized selection, with unadjusted 95\% paired mean-change intervals. Corrected signed-rank decisions are reported separately; a positive unadjusted mean interval is not a family-corrected penalty claim. Statistical families and differing model budgets are specified in Methods.}
\label{fig:task}\end{figure}

\subsection{Where does the electrode-count conclusion stop generalizing?}
The two split-derived fixed layouts gave different results. Odd-ID development/even-ID testing removed \FixedDropped, with change \FixedDiff\ (\FixedLCB\ to \FixedUCB), meeting the 0.02 criterion. The reverse removed \FixedSwapDropped, with change \FixedSwapDiff\ (\FixedSwapLCBSix\ to \FixedSwapUCB), whose lower bound missed the criterion. The reverse split met the percentile criterion in \FixedLayoutMonteCarloHits\ of \FixedLayoutMonteCarloN\ Monte Carlo repetitions (\FixedLayoutMonteCarloPercent\%), indicating numerical decision fragility rather than a probability that non-inferiority is true. These are two layouts rather than one layout validated twice; the count had already been chosen using the full cohort. Neither is an independent prospective validation.

Rest-inclusive, window-length and fixed-count random-forest sensitivities compare reduced systems with their own full arrays (table~\ref{tab:sensitivity}; figure~\ref{fig:limits}(B)). Complete window and rest curves appear in the Supplementary. They reuse the primary ranking and test robustness of reduced-model evaluation; they do not show that newly estimated rankings would be invariant to those changes. The random-forest check has no estimated minimum count because only twelve and sixteen were tested.

DB5 provided a small, separate-data test. Its count-12 lower bound under the percentile criterion was \DBFiveDeltaKTwelveLCBSix, close to the 0.02 boundary, and no count met the percentile-tail Holm rule at that margin. The exploratory null-centred Holm analysis likewise supported no reduced count, even at the looser 0.03 margin. With only \DBFiveSubjects\ participants and different hardware/tasks, it offers an uncertain boundary assessment rather than confirmation of the GRABMyo count. Derived bipolar pairs likewise failed to support a reduction from eight to seven at the primary margin. This representation cannot validate an equivalent reduction in physical bipolar sensors.

Full-array cross-session comparisons did not isolate placement displacement or reduced-layout robustness. Their participant results are retained in the Supplementary. A future independently collected same-montage cohort would be needed to validate a locked common layout, electrode count and online performance.

\begin{table}[htbp]
\centering\small
\caption{Within-condition reduced minus full-array comparisons at the 0.02 margin. Null denotes the exploratory null-centred studentized bootstrap with Holm correction across tested counts. The RF experiment tests only 12 and 16 electrodes, so it estimates no minimum count.}\label{tab:sensitivity}
\begin{adjustbox}{max width=\textwidth}
\begin{tabular}{lllllll}\toprule
Condition & $n$ & $k$ & $\bar\Delta$ & 95\% CI & $k^*_{\rm CI}$ & $k^*_{\rm null}$\\\midrule
Primary 250 ms & 43 & 12 & -0.0133 & [-0.0160, -0.0106] & 12 & 12\\
Rest included & 43 & 12 & -0.0139 & [-0.0166, -0.0111] & 12 & 12\\
150 ms & 43 & 12 & -0.0119 & [-0.0145, -0.0093] & 12 & 12\\
500 ms & 43 & 12 & -0.0126 & [-0.0161, -0.0091] & 12 & 12\\
NinaPro DB5 & 10 & 12 & -0.0126 & [-0.0200, -0.0056] & 12 & none\\
Derived bipolar & 43 & 7 & -0.0158 & [-0.0203, -0.0113] & none & none\\
Random forest & 43 & 12 & -0.0078 & [-0.0104, -0.0052] & not estimated & not estimated\\
\bottomrule\end{tabular}
\end{adjustbox}
\end{table}

\begin{figure}[htbp]\centering
\includegraphics[width=\textwidth]{../figures/vector/fig6_generalization_limits.pdf}
\caption{Boundaries of transfer. (A) Mean changes and intervals for two distinct split-derived layouts, using 21 and 22 test participants and all development windows for fitting. (B) Twelve-electrode changes for task, window and classifier sensitivities, with each condition's own sixteen-electrode reference. (C) DB5 reduction versus its sixteen-channel reference ($n=10$). (D) Derived bipolar reduction versus its eight-pair reference ($n=43$). The dashed line is the 0.02 margin. Every interval resamples participants; scores are not pooled across hardware or representations.}
\label{fig:limits}\end{figure}

\section{Discussion}
The study supports removing one quarter of the GRABMyo monopolar recording sites within a defined mean recognition tolerance under this offline LOSO protocol. It does not show that every participant or gesture is preserved. The observed individual losses and corrected gesture declines make this distinction practically relevant: count selection for a real interface should include task priorities and user-specific failure criteria rather than relying solely on a pooled mean.

Attribution supplied predictive information beyond the mean random-subset baseline, and deletion with retraining supported its ranking relevance. Spatial coverage was a more competitive comparator. The matched primary experiment places geometric12's mean loss near the tolerance boundary; the direct mean advantage of SHAP12 remains uncertain. The full geometric curve required \RevGeomKStarCI\ retained sites under the primary margin and decision rule, but only one deterministic geometric ordering was tested. These results motivate attribution as one evaluated selection procedure, with additional computational cost, rather than establishing a universal advantage over spatial selection. The shared-model experiment found a positive SHAP mean advantage over grouped permutation at six retained electrodes, but no established mean advantage at twelve. The gain under more aggressive reduction does not mean that six electrodes met the primary tolerance. It changes permutation grouping, participant-disjoint validation, scaler isolation and resources together, so it cannot attribute differences solely to computational budget.

Stability regularization was not an established improvement. Five-fold pooling obscures some person heterogeneity, which motivated the participant-level extension, but more descriptive dispersion does not imply that subtracting it improves classification. All penalty settings were fixed and disclosed rather than selected on outer scores. Correlated electrode features, normalization and shared attributing models complicate physiological interpretation; high attribution is model reliance and can reflect redundancy, task discrimination or montage geometry without identifying a particular muscle.

The inference has two distinct limits. First, the tolerance is an engineering choice, not a clinically validated equivalence bound; tighter or looser tolerances support different counts. Second, participant bootstrap intervals are conditional on fitted models with strongly overlapping training sets. The null-centred sensitivity corrects how the bootstrap test distribution is constructed, but cannot remove training-set dependence or development exposure. A nested development cohort, independently locked tolerance and external validation would strengthen future inference.

Several aspects restrict transfer. Only able-bodied public datasets were analyzed, not target clinical users. Zero-phase filtering is non-causal. Complete rest/window curves retain the primary ranking; the classifier sensitivity also retains it and samples only two counts. DB5 uses different sensors, sampling, labels and a small cohort, and derived bipolar pairs are constructed signals. The two split layouts differ and count selection had already seen the cohort. Full-array cross-session results do not demonstrate reduced-layout displacement robustness. No real-time accuracy, latency, energy, hardware cost reduction or clinical benefit was measured. Clinical channel-count experiments provide a further reason to keep this scope narrow: real-time prosthesis control can benefit from additional channels even when offline reduction appears promising \citep{simon2024implications}. Such evidence uses a different cohort, interface and endpoint and cannot validate or refute the present offline mean tolerance.

A prospective follow-up should fix a single geometry and count using development data only, record new participants with the documented montage, and define participant- and gesture-level loss criteria before testing. It should evaluate causal preprocessing and a closed-loop task, recording prediction latency and electrode failures as measured outcomes.

\section{Conclusion}
Cross-fitted attribution supported reduction of the GRABMyo monopolar array to \KStarMarginPtTwo\ electrodes at a 0.02 mean Macro-F1 tolerance under participant-independent offline evaluation. Utility beyond random selection was supported. Under shared resources, SHAP exceeded grouped permutation at six electrodes; its twelve-electrode mean advantage remained uncertain, as did the primary SHAP12--geometry12 contrast. Stability regularization did not establish an improvement. Participant losses, gesture-specific degradation, distinct fixed layouts and representation-dependent results restrict the claim to an evaluated reduction procedure. Independent validation is still required for a locked physical interface.

\data{Public GRABMyo v1.1.0 recordings are available from PhysioNet (\url{https://doi.org/10.13026/89dm-f662}); NinaPro DB5 recordings are available from \url{https://ninapro.hevs.ch}.}
\suppdata{The Supplementary contains signal-processing and inference details, exploratory analyses, complete result tables and additional figures.}

\begingroup\small\setlength{\bibsep}{2pt}\bibliographystyle{plainnat}\bibliography{references}\endgroup
\end{document}
